% SOURCE_AUTHORITY: sga5_fr_workpass.tex, lines 11943--11951 (LF-normalized unit).
% SOURCE_SHA256: 791F4EFFC5E02832D5D77ED03518C8156D6F07E4C8238B03545DB93D883FBB28
\subsection*{6. Complexos perfeitos}

Um objeto da categoria $D(\cC)$ será chamado complexo $\cC_o$-perfeito se for isomorfo a um objeto de $D_o^b(\cC)$, isto é, a um complexo limitado cujas componentes estão todas em $\cC_o$. Denotaremos por $D_o(\cC)_{\parf}$ a subcategoria plena de $D(\cC)$ gerada pelos complexos perfeitos.

A categoria $D_o(\cC)_{\parf}$ não é, portanto, mais do que a imagem essencial de $D_o^b(\cC)$ pelo funtor de inclusão
\[
D_o^b(\cC)\longrightarrow D(\cC).
\]
É fácil ver que $D_o(\cC)_{\parf}$ é uma subcategoria triangulada de $D(\cC)$ e que, portanto, possui uma estrutura de categoria triangulada.
